\documentclass{article}
\usepackage[utf8]{inputenc}

\title{The Elements of Typographic Style}
\author{Robert Bringhurst}
\date{1992}

\begin{document}
\maketitle

\pgmark{1}\section*{Introduction}

Typography exists to honor content. The page is a stage for the words; the
margins are the wings; the leading is the air through which the reader
breathes. This brief excerpt --- one of several reproduced here for the
purpose of exercising the Virgil Library's rendering surface --- is taken
from the indexed copy stored in \texttt{papers/bringhurst1992/}.

\pgmark{2}\section{The grand design}

A well-made book is a small architecture\footnote{Compare Tschichold's
\emph{Form of the Book}, where the same metaphor recurs.} in which type,
margin, and leading conspire to make the text comfortable for the eye.

\pgmark{3}When the body of the text is set in a serif of moderate
weight --- $11/14$~pt is a common choice --- and the measure is held to
between $60$ and $75$ characters per line, the page can be read for hours
without fatigue. \cite{genette1997} observed something analogous about
paratexts.

\pgmark{4}\section{Margins}

The classical canons of page construction --- van de Graaf, the Villard
diagram, the golden section --- all begin from the premise that the
text block should be smaller than the leaf, and that the inner margin
should be smaller than the outer.\footnote{This is also the rationale
for the \emph{recto-verso asymmetry} that troubles modern reflowable
formats.}

\pgmark{5}A page that has been crowded to its edges has nothing to give
the reader. A page that has breathing room around the text invites the
reader in.

\pgmark{6}\section{The footnote as a quiet companion}

A footnote is the shadow that walks beside a sentence. It belongs to the
sentence, but it does not interrupt it. \cite{grafton2020}'s history is
the standard treatment, though the typography is older than the genre.

\pgmark{7}When a footnote grows past a quarter of the page, the typesetter
faces a choice: hyphenate the body to make room, push the footnote
forward to the next page, or convert the apparatus to endnotes.

\pgmark{8}\section{Closing}

This file is a stub for the Virgil Library dev preview --- enough text
to exercise paragraph rendering, the \verb|\pgmark{}| margin chips,
citations, footnotes, and headings, but no more.

\end{document}
